% Draft 2026-10-17 (writer), revised after Codex review and to the researcher's style (notes/style_guide.md).
% No results yet: the evaluation is described as a protocol. Once results are frozen, add one sentence with
% \resSLAViolQuantum, \resSLAViolDInos, \resParamsQuantum, \resParamsDInos, \resParamsControl and \resNumSeeds.
Autoscaling a Redis cluster on Kubernetes is a control problem with delays. A new pod helps only after it has
started and joined the cluster. A threshold rule that reacts to the current CPU level, such as the Horizontal
Pod Autoscaler (HPA) scaling on CPU, therefore tends to add capacity after the load has already risen. Deep
Reinforcement Learning (Deep RL) agents such as DInos learn proactive policies from a window of recent metrics
and adapt them to a new environment by fine-tuning. In the DInos implementation, almost all of the Q-network's
parameters sit in a recurrent encoder, and fine-tuning updates all of them.
In this paper, we introduce Sofia, a hybrid quantum--classical Deep RL agent in which a variational quantum
circuit replaces part of the DInos Q-network.\todo{Name the quantum component once D5 (candidate selection) is decided.}
To evaluate it, we build a simulator in which latency follows from a queueing model of the Redis pods, and pod
start-up and data rebalancing take time. Before any agent is compared, the simulator must pass sanity gates:
trivial fixed policies have to lose to a tuned threshold autoscaler.
Our protocol compares Sofia with DInos retrained in the same simulator, a threshold policy modeled on the HPA,
and a classical network with as many trainable parameters as Sofia. It measures SLA violations, tail latency,
sample efficiency, adaptation and model size, and checks whether the ranking holds on a real
cluster.\todo{Results sentence with macros once results are frozen.}
All circuits run on a classical simulator, and we make no claim about quantum speed-up.
